%%%%%%%%%%%%%%%%%%%%%%%%%%%%%%%%%%%%%%%%%%%%%%%%%%%%%%%%%%%%%%%%%%%%%%%%%%%%%%%%
%%
\cleardoubleoddpage%  Make sure to start each chapter on a new odd page
\chapter{PyTorch Reimplementation}
\label{sec:pytorchreimplementation}

\section{Design of the PyTorch Version}

The PyTorch implementation \cite{paszke2019pytorch} is organized around
three components: a plain \texttt{torch.nn.Module} multilayer perceptron
mapping a parameter vector
$x \in [-1,1]^d$ to a flat vector of finite element degrees of freedom
(Section~\ref{sec:methodology}), a training loop implementing the paper's
optimizer, learning-rate schedule, and checkpoint/restore rule, and a
sweep driver that iterates over the eight PDE cases, six architectures,
fourteen training-set sizes, and twelve trials, writing one metrics row
per run so that a sweep can be resumed after an interruption without
recomputing already-completed runs. Training and test data are read
directly from the \texttt{.npz} files and mass-matrix sidecars produced by
the FEniCS/sparse-grid data generation pipeline (Section~\ref{sec:methodology});
the PyTorch process itself never depends on FEniCS or Tasmanian, keeping
the legacy-library constraints confined to the WSL2 data-generation
environment.

\section{Key Implementation Decisions}

Two implementation choices required deviating from PyTorch's defaults to
match the paper's stated training protocol rather than a generic one.
First, weight initialization: PyTorch's built-in Kaiming initialization
selects its gain based on the target activation function, whereas Keras'
\texttt{HeUniform} initializer -- the initializer implicitly used by the
paper's TensorFlow/Keras reference implementation -- always uses the
ReLU-derived gain, regardless of the network's actual activation. The
PyTorch implementation therefore fixes the nonlinearity argument passed to
Kaiming initialization to \texttt{"relu"} uniformly across all three
activation functions used in the architecture sweep (ReLU, ELU, tanh), so
that the initial weight distribution matches the paper's rather than
PyTorch's own convention. Second, per-trial seeding: each of the twelve
trials per case must be reproducible individually and distinct from the
other eleven, so a \texttt{trial\_seed} argument is threaded through the
sweep driver into both the training-subsample draw and
\texttt{torch.manual\_seed} before weight initialization, rather than
seeding the whole sweep once at startup.

The paper's checkpoint/restore rule (Section~\ref{sec:methodology}) is
implemented as an explicit two-trigger condition evaluated once per epoch:
the current \texttt{state\_dict} is deep-copied into a checkpoint buffer
whenever the loss achieves a new best value \emph{or} the ratio between the
current loss and the loss at the last checkpoint drops below $1/8$: after
the epoch loop terminates (either by reaching the epoch cap or by falling
below the loss tolerance), the checkpointed weights are loaded back into
the model if the final-epoch loss is worse than the checkpoint's loss.

\section{Deviations from the Original}

Beyond the deep-learning framework substitution shared with the JAX
implementation (Chapter~\ref{sec:jaximplementation}), a small number of
deviations are specific to how the PyTorch implementation interacts with
the FEniCS-generated data and are documented in detail in
Appendix~\ref{app:an-appendix}: the diffusion mesh is a structured mesh
rather than the unstructured mesh the paper's figures suggest, because the
unstructured mesh generator \texttt{mshr} is not usable in the installed
FEniCS environment; the NSB inlet velocity profile substitutes a
coordinate reflection of the paper's printed formula, which is otherwise
degenerate at the inlet boundary; and the NSB pressure field, which the
paper's formulation only determines up to an additive constant via a
Lagrange multiplier, is instead fixed by a post-hoc mean-shift enforcing
$\int_\Omega p = 0$ exactly. None of these deviations depend on the choice
of deep learning framework -- they occur upstream, in data generation --
so they apply identically to the JAX implementation's training data.

\section{Verification of Correctness}

The PyTorch implementation served as the first target for verifying the
full pipeline end to end, since it is easier to step through and debug
than a JIT-compiled JAX equivalent. Verification proceeded in stages:
unit tests check the initialization-distribution bounds
($\sqrt{6/\text{fan\_in}}$ across all three activations), the
checkpoint/restore rule's two triggers individually and jointly on
synthetic non-monotonic loss curves, and per-trial seed determinism and
distinctness. At the PDE level, single finite element solves were compared
visually against the paper's own figures (the diffusion solution profile
against Fig.~5, the NSB circulation pattern against Fig.~6) before any
dataset-scale generation was run, to catch a wrong element family or
boundary condition as early and cheaply as possible. The resulting PyTorch
metrics also serve as the correctness reference against which the JAX
implementation's numerical agreement is assessed
(Section~\ref{sec:jaximplementation}).
%%
%%%%%%%%%%%%%%%%%%%%%%%%%%%%%%%%%%%%%%%%%%%%%%%%%%%%%%%%%%%%%%%%%%%%%%%%%%%%%%%%
